\documentclass[UTF8,11pt]{ctexart}
\usepackage{amsmath,amssymb,geometry}
\geometry{margin=22mm}
\begin{document}
\section*{P4: 论文主分析使用的鲁棒上界近似模型（给定运营规模参数）}
\paragraph{来源与范围}依据原论文PDF第6, 7, 8, 9, 10, 11, 12, 13, 14, 15, 16, 17, 18页，第2节式(7)、(8)–(13)用于解释来源，保留第2.4节式(14)。

\paragraph{所保留形式}
保留原文第2.4节式(14)：给定批量$q$及拟参与司机数$N$时，选择众包比例和忙时工资的鲁棒上界规划。
该模型是原文式(7)的鲁棒上界近似。
\subsection*{输入与必要定义}
$m$为配送批次数，$q$为每批订单数，$N$为司机数，$m=N\tau$，$m,q,N,\tau\in\mathbb Z_{++}$。
第$j$批客户为$B_j=\{(j-1)q+1,\ldots,jq\}$，配送速度$v_j>0$。
订单到达间隔与现场服务时间的均值和标准差为$\mu_i^a>0,\sigma_i^a\ge0,\mu_i^s\ge0,\sigma_i^s\ge0$。
$\beta\ge0$为论文将地理密度及TSP上界合并后的长度系数，$t_j=\beta\sqrt{q+1}/v_j$。
$\alpha>0$为承诺配送期限，$0<\epsilon<1$为利用率安全余量，$c_T\ge0$为第三方每单配送费。
$\gamma^a,\gamma^s,\gamma^k\ge0$为到达、服务及机会成本的保守程度；
$\mu_{(r)}^k$为司机机会成本均值从小到大排列的前$N$项，$\sigma^k\ge0$为共同标准差。
司机$h$的服务批次索引$J_h=\{j_{h1},\ldots,j_{h\tau}\}$给定，每个$J_h$从每轮$N$批中取一批，且各$J_h$划分全部$m$批；令$D_h=\bigcup_{j\in J_h}B_j$。
定义固定系数
\[
 C=\frac{\sum_{j=1}^m t_j+\sum_{h=1}^N\left(\sum_{i\in D_h}\mu_i^s+\gamma^s\sqrt{\sum_{i\in D_h}(\sigma_i^s)^2}\right)}{mq},\quad
 R=\frac{\sum_{i=1}^{mq}\mu_i^s+\sum_{j=1}^m t_j}{N\sum_{i=1}^{mq}\mu_i^a},\quad
 K_N=\frac{\sum_{r=1}^N\mu_{(r)}^k+\gamma^k\sqrt N\sigma^k}{N}.
\]
对$j>N$，$a_j=\lfloor(j-1)/N\rfloor$。对$1\le i\le a_j$，
$A_{ji}=\{(j-iN)q+1,\ldots,(j-1)q+1\}$，
$D_{ji}=\bigcup_{\ell=0}^{i}B_{j-\ell N}$。
服务时间上界由以下表达式定义。
\begin{align*}
G_j(P)&=\frac{\sum_{r=(j-1)q+2}^{jq}\mu_r^a+\gamma^a\sqrt{\sum_{r=(j-1)q+2}^{jq}(\sigma_r^a)^2}}{P}
+\sum_{r\in B_j}\mu_r^s+\gamma^s\sqrt{\sum_{r\in B_j}(\sigma_r^s)^2}+t_j,\\
F_{ji}(P)&=\frac{-\sum_{r\in A_{ji}}\mu_r^a+\gamma^a\sqrt{\sum_{r\in A_{ji}}(\sigma_r^a)^2}}{P}
+\sum_{r\in D_{ji}}\mu_r^s+\gamma^s\sqrt{\sum_{r\in D_{ji}}(\sigma_r^s)^2}+\sum_{\ell=0}^{i}t_{j-\ell N}.
\end{align*}
空求和为0，特别是$q=1$时的等候组团项。
\subsection*{唯一核心规划}
决策$P$为众包订单比例，$w$为司机每单位忙时工资。
\begin{align}
\min_{P,w}\quad &CPw+(1-P)c_T,\tag{14a}\\
\text{s.t.}\quad&G_j(P)\le\alpha &&j=1,\ldots,m,\tag{14b}\\
&F_{ji}(P)\le\alpha &&j=N+1,\ldots,m,\quad i=1,\ldots,a_j,\tag{14c}\\
&RP\le1-\epsilon,\tag{14d}\\
&RPw\ge K_N,\tag{14e}\\
&0<P\le1,\quad w>0.\tag{14f}
\end{align}
逐项$F_{ji}\le\alpha$与原式对$i$取最大值后约束等价。
$q,N$是给定的运营规模参数。

\end{document}
